\section{What the corner sees}\label{sec:corner}

At a corner, the two dimensional interpolation kernel simplifies in a
useful way: its absolute values become a positive Cauchy sum.  There are
still two adjacent meshes, of sizes $n$ and $n+1$.  Keeping that small
difference, and the endpoint weights, will be essential for the first
asymptotic correction.

Throughout this section let $n\geq2$, $m=n+1$, and write
\[
 \theta_j=\frac{j\pi}{n},\qquad \phi_k=\frac{k\pi}{m},\qquad
 \varepsilon_j=\begin{cases}\tfrac12,&j=0,n,\\1,&0<j<n,\end{cases}
 \quad
 \eta_k=\begin{cases}\tfrac12,&k=0,m,\\1,&0<k<m.\end{cases}
\]
The node set uses $j+k$ even.  We denote the common value at the two
corners outside this set by $\Lambda_c(n)$; the argument below also
verifies that the two values agree.  In particular $(-1,-1)$ is such a
corner for every $n$.

\begin{proposition}[Positive corner formula]\label{prop:corner-formula}
Put
\[
 f(u,v)=\frac{1+\cos v}{|\cos u-\cos v|},\qquad u\ne v.
\]
Then
\begin{equation}\label{eq:corner-positive-sum}
 \Lambda_c(n)=\frac4{nm}
 \sum_{\substack{0\leq j,k\leq n\\j+k\ \mathrm{odd}}}
 \varepsilon_j\eta_k f(\theta_j,\phi_k)+\frac1m.
\end{equation}
Every denominator in this sum is nonzero.  The other two corner values
are $1$.
\end{proposition}

\begin{proof}
We first sum the corner kernel in \eqref{eq:cardinal-kernel}.  Define
\[
 D_n(u,v)=\sum_{|p|+|q|\leq n}e^{i(pu+qv)}.
\]
Grouping the four sign choices shows that this is the product
Chebyshev sum occurring at a corner in that formula.  Set
$S=(u+v)/2$, $T=(v-u)/2$.  The substitution
$(p,q)\mapsto(p+q,q-p)$ transforms the diamond into the two squares
with equal coordinate parity.  Finite geometric sums therefore give
\begin{equation}\label{eq:corner-rhombus-product}
 D_n(u,v)=
 \frac{\sin((n+1)S)\sin((n+1)T)+\sin(nS)\sin(nT)}
 {\sin S\sin T}.
\end{equation}
Each quotient here is its continuous trigonometric polynomial at a
vanishing denominator.

Reflect the node indices about a corner which is not a node.  The
reflected indices form the odd checkerboard, the endpoint weights are
unchanged, and \eqref{eq:cardinal-kernel} becomes
\[
 \sum_{j+k\ \mathrm{odd}}\frac{2\varepsilon_j\eta_k}{nm}
       |D_n(\theta_j,\phi_k)-(-1)^j|.
\]
For $j+k$ odd and $u=\theta_j$, $v=\phi_k$ with $u\ne v$,
product to sum in \eqref{eq:corner-rhombus-product} yields
\[
 D_n(u,v)=(-1)^j\frac{2+\cos u+\cos v}{\cos u-\cos v},
 \qquad |D_n(u,v)-(-1)^j|=2f(u,v).
\]
The meshes can coincide only at $(j,k)=(0,0)$ or $(n,m)$, since
$n$ and $m$ are coprime.  The odd checkerboard excludes the first.
At the second the continuous value is
$D_n(\pi,\pi)=(-1)^n(2n+1)$, so its corrected absolute value is
$2n$.  Its weight is $1/(2nm)$, giving the separate $1/m$ in
\eqref{eq:corner-positive-sum}.  For $k=m$ and $j\ne n$ the numerator
of $f$ is zero.  This proves the formula and its independence of the
chosen nonnode corner.  Cardinality gives value $1$ at a node corner.
\end{proof}

\subsection{The part of each row that matters}

An absolute row can be understood by finding where its sign changes.
This is particularly simple here, and gives a useful alternative to
the double sum.  For $0<x<1$ put
\begin{equation}\label{eq:corner-w-B}
 w(x)=\cot\frac{\pi x}{2},\qquad
 B(x)=\cot(\pi x)-\frac1{\pi x}+\frac1{\pi(1-x)}.
\end{equation}
Both endpoint poles have been removed from $B$: it is analytic on a
neighborhood of $[0,1]$, with $B(0)=1/\pi$ and $B(1)=-1/\pi$.
For $x_r=r/m$ and $\alpha_r=(1-x_r)/2$ define
\begin{align}
 T_{n,r}&=\sum_{\ell=0}^{r-1}
             \cot\frac{\pi(\ell+\alpha_r)}n,
 &V_n&=\sum_{r=1}^n w(x_r)
             \sum_{\ell=0}^{r-1}B\!\left(\frac{\ell+\alpha_r}{n}\right).
 \label{eq:corner-partial-rows}
\end{align}
The upper limit is $r-1$, rather than $n-1$: these are precisely the
positive samples of the row.

We use the harmonic series definition
\begin{equation}\label{eq:corner-psi-definition}
 \psi(z)=-\gamma+\sum_{k=0}^{\infty}
       \left(\frac1{k+1}-\frac1{k+z}\right),\qquad z>0,
\end{equation}
where $\gamma=\lim_{k\to\infty}(\sum_{j=1}^k j^{-1}-\log k)$.
Set, for $s=m-r$,
\begin{align}
 D_{m,r}={}&\psi(r+\alpha_r)+\psi(s-\alpha_r)
                  -\psi(m-\alpha_r)-\psi(\alpha_r),
 \label{eq:corner-harmonic-row}\\
 \mathcal H_m={}&\frac4{\pi m}\sum_{r=1}^{m-1}w(r/m)D_{m,r}.
 \label{eq:corner-harmonic-triangle}
\end{align}

\begin{proposition}[Exact row decomposition]
\label{prop:corner-harmonic-decomposition}
For every $n\geq2$,
\begin{align}
 \Lambda_c(n)&=1-\frac{2n}{m}
             +\frac4{nm}\sum_{r=1}^n w(r/m)T_{n,r},
             \label{eq:corner-cotangent-decomposition}\\
 &= -1+\frac2m+\mathcal H_m+\frac{4V_n}{nm}.
             \label{eq:corner-exact-decomposition}
\end{align}
These identities retain all rows, including $r=n$, and both endpoint
contributions.
\end{proposition}

\begin{proof}
Fix $1\leq r\leq n$, put $v=r\pi/m$, and extend the first mesh
periodically.  Its $n$ samples have $u_j=j\pi/n$, with
$j\equiv1-r\pmod2$ modulo $2n$.  Write
\[
 F_r(u)=\frac{1+\cos v}{\cos u-\cos v},\qquad
 R_r=\sum_j|F_r(u_j)|.
\]
There is no collision.  The identity
\[
 F_r(u)=\frac{\cot(v/2)}2
       \left(\cot\frac{v-u}{2}+\cot\frac{v+u}{2}\right)
\]
and the finite cotangent sum
$\sum_{\ell=0}^{n-1}\cot(z+\ell\pi/n)=n\cot(nz)$ give
$\sum_jF_r(u_j)=n$.  The cotangent sum follows by taking the
logarithmic derivative of the sine product over the $n$ equally
spaced zeros; alternatively it follows from a finite geometric sum.
In the present substitution its right hand cotangent is $\tan(v/2)$,
which cancels $\cot(v/2)$.

The positive samples have centered representatives $|j|<nr/m$.
Since $r-1<nr/m<r$, they are exactly
\[
 j=-r+1,-r+3,\ldots,r-1.
\]
If their sum is $P_r$, the signed row identity gives $R_r=2P_r-n$.
Their symmetry in $u$ and the displayed identity for $F_r$ give
$P_r=w(r/m)T_{n,r}$, by writing $j=r-1-2\ell$.

Folding this periodic row back to $0\leq j\leq n$ gives half of
$R_r$, with exactly the weights $\varepsilon_j$.  The row $r=0$
has $R_0=n^2$, from
$\sum_{\ell=0}^{n-1}\csc^2((\ell+1/2)\pi/n)=n^2$ (differentiate
the finite cotangent identity).  Its second-coordinate weight is
$1/2$.  Together with the exceptional $1/m$ in
\eqref{eq:corner-positive-sum}, this gives
\[
 \Lambda_c(n)=\frac1m+\frac nm+\frac2{nm}\sum_{r=1}^n(2P_r-n),
\]
which is \eqref{eq:corner-cotangent-decomposition}.

Finally subtract the two poles of each cotangent in $T_{n,r}$.
The recurrence from \eqref{eq:corner-psi-definition} gives exactly
\[
 T_{n,r}=\frac n\pi D_{m,r}
       +\sum_{\ell=0}^{r-1}B\!\left(\frac{\ell+\alpha_r}{n}\right).
\]
The scalar term $1-2n/m$ equals $-1+2/m$, proving
\eqref{eq:corner-exact-decomposition}.
\end{proof}

\subsection{A one dimensional reference and the limiting defect}

If the two meshes in \eqref{eq:corner-positive-sum} were identical,
the answer would be a square.  Define
\begin{equation}\label{eq:corner-L-definition}
 L_n=\frac2n\sum_{\substack{1\leq r\leq n\\r\ \mathrm{odd}}}
       d_r\csc\frac{r\pi}{2n},\qquad
 d_r=\begin{cases}\tfrac12,&r=n,\\1,&r<n.\end{cases}
\end{equation}
The half weight matters when $n$ is odd.  The following facts identify
the constant which remains after the leading square is removed.  Their
proofs, including the endpoint estimates, are given in
Appendix~\ref{sec:asym-details-reference}.

\begin{lemma}[The reference square]\label{lem:corner-reference-square}
For every $n\geq2$,
\[
 \frac4{n^2}\sum_{\substack{0\leq j,k\leq n\\j+k\ \mathrm{odd}}}
      \varepsilon_j\varepsilon_k f(\theta_j,\theta_k)=L_n^2.
\]
\end{lemma}

\begin{lemma}[Size of the reference]\label{lem:corner-one-dimensional}
Let $a=2/\pi$ and $d=\gamma+\log(8/\pi)$.  On the full integer
sequence $n\geq1$,
\begin{equation}\label{eq:corner-L-estimate}
 \left|L_n-a(\log n+d)\right|\leq\frac1{4n^2}.
\end{equation}
\end{lemma}

\begin{lemma}[The defect constant]\label{lem:corner-defect-limit}
As $n\to\infty$ through all integers,
\begin{equation}\label{eq:corner-kappa-limit}
 \Lambda_c(n)-L_n^2\longrightarrow\kappa,
 \qquad
 \kappa=\frac8\pi\int_0^1 t^2\cot\frac{\pi t}{2}
       \sum_{\substack{\ell\geq1\\\ell\ \mathrm{odd}}}
              \frac1{\ell(\ell^2-t^2)}\,dt.
\end{equation}
The integrand has continuous limits at $0$ and $1$.
\end{lemma}

The mechanism behind \eqref{eq:corner-kappa-limit} is easy to see
before making the estimates.  Fix the odd distance $\ell=|j-k|$
from the diagonal and put $t=j/n$.  The change of mesh moves the
two neighboring denominators from $\ell$ to $\ell-t$ and
$\ell+t$.  Pairing the two directions leaves
\[
 \frac1{\ell-t}+\frac1{\ell+t}-\frac2\ell
       =\frac{2t^2}{\ell(\ell^2-t^2)}.
\]
This is precisely the positive defect in the integral.  Its apparent
pole at $\ell=t=1$ is canceled by the zero of
$\cot(\pi t/2)$.  Thus the limiting constant comes from the
adjacent meshes, while the leading logarithms come from the square.
